El fraude financiero digital constituye uno de los principales riesgos asociados al uso masivo de tecnologías y redes sociales, tanto por su impacto económico como por su capacidad de adaptación a nuevos entornos digitales. En los últimos años, plataformas como X (antes Twitter) no solo han actuado como canal para la ejecución de estafas, sino también como espacio donde los usuarios comparten experiencias, denuncias y advertencias sobre distintos tipos de fraude, generando un corpus textual de enorme valor para la vigilancia temprana.

Este Trabajo Fin de Máster propone un enfoque metodológico basado en Procesamiento del Lenguaje Natural (\textit{Natural Language Processing}, NLP) para identificar patrones discursivos asociados al fraude financiero a partir de publicaciones en redes sociales. El objetivo es detectar señales tempranas y tipologías emergentes de estafa, y evaluar la transferibilidad de los patrones observados en el contexto estadounidense al contexto español.

Para ello se utilizan datos extraídos mediante la API de X correspondientes a Estados Unidos durante el año 2025, país que actúa como fuente de aprendizaje por la madurez de su ecosistema digital. A la luz de los 122.223 incidentes gestionados por INCIBE en 2025, de los cuales cuatro de cada diez corresponden a fraude digital, España opera como contexto de destino en el que se evalúa la transferibilidad del modelo, igualmente para esto se usan datos extraídos mediante la API de X. Sobre el corpus estadounidense se aplica un pipeline que combina limpieza textual, \textit{topic modeling} no supervisado mediante BERTopic, clasificación automática en categorías de fraude mediante un modelo \textit{zero-shot} basado en BART-MNLI, y traducción automática al español mediante modelos Marian (Helsinki-NLP). Los resultados se analizan en términos de frecuencia, co-ocurrencia léxica y evolución temporal de las tipologías detectadas.

Los hallazgos muestran la existencia de estructuras discursivas estables en las menciones de fraude financiero en redes sociales, independientes de la plataforma o el contexto geográfico, lo que sugiere la viabilidad de trasladar indicadores lingüísticos entre jurisdicciones. Este enfoque demuestra el potencial de las redes sociales como sensor temprano de riesgos financieros y aporta una metodología complementaria a los sistemas institucionales de alerta basada en el análisis del lenguaje social.

\textbf{Palabras Clave:} fraude financiero, procesamiento del lenguaje natural, BERTopic, clasificación \textit{zero-shot}, redes sociales, X (Twitter), transferencia translingüística, alerta temprana, España, Estados Unidos.
